\chapter{Calculus Review}
\label{ch:calculus}
% ============================================================================

\section{Differential Calculus}

\subsection{Derivatives and Rules}

The derivative of a function $f(t)$ represents its instantaneous rate of change:
\begin{equation}
\frac{\dd f}{\dd t} = f'(t) = \lim_{\Delta t \to 0} \frac{f(t + \Delta t) - f(t)}{\Delta t}
\end{equation}

\subsubsection{Key Differentiation Rules}

\begin{table}[htbp]
\centering
\begin{tabular}{ll}
\toprule
\textbf{Rule} & \textbf{Formula} \\
\midrule
Power Rule & $\frac{\dd}{\dd t}(t^n) = n t^{n-1}$ \\
Exponential & $\frac{\dd}{\dd t}(\ee^{at}) = a\ee^{at}$ \\
Logarithm & $\frac{\dd}{\dd t}(\ln t) = \frac{1}{t}$ \\
Sine & $\frac{\dd}{\dd t}(\sin\omega t) = \omega\cos\omega t$ \\
Cosine & $\frac{\dd}{\dd t}(\cos\omega t) = -\omega\sin\omega t$ \\
Product Rule & $\frac{\dd}{\dd t}(fg) = f'g + fg'$ \\
Quotient Rule & $\frac{\dd}{\dd t}\left(\frac{f}{g}\right) = \frac{f'g - fg'}{g^2}$ \\
Chain Rule & $\frac{\dd}{\dd t}f(g(t)) = f'(g(t)) \cdot g'(t)$ \\
\bottomrule
\end{tabular}
\caption{Essential differentiation rules for control systems.}
\label{tab:diff_rules}
\end{table}

\subsection{Partial Derivatives}

For functions of multiple variables $f(x_1, x_2, \ldots, x_n)$, partial derivatives hold all variables constant except one:
\begin{equation}
\frac{\partial f}{\partial x_i} = \lim_{\Delta x_i \to 0} \frac{f(x_1, \ldots, x_i + \Delta x_i, \ldots, x_n) - f(x_1, \ldots, x_n)}{\Delta x_i}
\end{equation}

The \textbf{gradient} is the vector of all partial derivatives:
\begin{equation}
\nabla f = \begin{bmatrix} \frac{\partial f}{\partial x_1} & \frac{\partial f}{\partial x_2} & \cdots & \frac{\partial f}{\partial x_n} \end{bmatrix}^\top
\end{equation}

The \textbf{Hessian} is the matrix of second partial derivatives:
\begin{equation}
\mH = \begin{bmatrix}
\frac{\partial^2 f}{\partial x_1^2} & \frac{\partial^2 f}{\partial x_1 \partial x_2} & \cdots \\
\frac{\partial^2 f}{\partial x_2 \partial x_1} & \frac{\partial^2 f}{\partial x_2^2} & \cdots \\
\vdots & \vdots & \ddots
\end{bmatrix}
\end{equation}

\subsection{Taylor Series Expansion}

Taylor series are essential for linearization of nonlinear systems:
\begin{equation}
f(x) = f(a) + f'(a)(x-a) + \frac{f''(a)}{2!}(x-a)^2 + \frac{f'''(a)}{3!}(x-a)^3 + \cdots
\end{equation}

For multivariate functions about point $\vecx_0$:
\begin{equation}
f(\vecx) \approx f(\vecx_0) + \nabla f(\vecx_0)^\top (\vecx - \vecx_0) + \frac{1}{2}(\vecx - \vecx_0)^\top \mH(\vecx_0) (\vecx - \vecx_0)
\end{equation}

\begin{keyconceptbox}[title=Linearization in Control]
For nonlinear systems $\dot{\vecx} = \vecf(\vecx, \vecu)$, we linearize about an equilibrium point $(\vecx_0, \vecu_0)$:
\begin{equation}
\Delta\dot{\vecx} \approx \mA \Delta\vecx + \mB \Delta\vecu
\end{equation}
where $\mA = \frac{\partial \vecf}{\partial \vecx}\big|_{(\vecx_0, \vecu_0)}$ and $\mB = \frac{\partial \vecf}{\partial \vecu}\big|_{(\vecx_0, \vecu_0)}$
\end{keyconceptbox}

\section{Integral Calculus}

\subsection{Definite and Indefinite Integrals}

The indefinite integral (antiderivative):
\begin{equation}
\int f(t) \, \dd t = F(t) + C \quad \text{where} \quad F'(t) = f(t)
\end{equation}

The definite integral:
\begin{equation}
\int_a^b f(t) \, \dd t = F(b) - F(a)
\end{equation}

\subsubsection{Key Integration Formulas}

\begin{table}[htbp]
\centering
\begin{tabular}{ll}
\toprule
\textbf{Function} & \textbf{Integral} \\
\midrule
$t^n$ $(n \neq -1)$ & $\frac{t^{n+1}}{n+1} + C$ \\
$\frac{1}{t}$ & $\ln|t| + C$ \\
$\ee^{at}$ & $\frac{1}{a}\ee^{at} + C$ \\
$\sin\omega t$ & $-\frac{1}{\omega}\cos\omega t + C$ \\
$\cos\omega t$ & $\frac{1}{\omega}\sin\omega t + C$ \\
$\frac{1}{a^2 + t^2}$ & $\frac{1}{a}\arctan\frac{t}{a} + C$ \\
\bottomrule
\end{tabular}
\caption{Essential integration formulas.}
\label{tab:int_rules}
\end{table}

\subsection{Integration by Parts}

\begin{equation}
\int u \, \dd v = uv - \int v \, \dd u
\end{equation}

This technique is essential for deriving Laplace transform properties.

\subsection{Convolution Integral}

The convolution of two functions $f(t)$ and $g(t)$ is:
\begin{equation}
(f \conv g)(t) = \int_{-\infty}^{\infty} f(\tau) g(t - \tau) \, \dd\tau
\end{equation}

For causal systems (signals zero for $t < 0$):
\begin{equation}
(f \conv g)(t) = \int_0^t f(\tau) g(t - \tau) \, \dd\tau
\end{equation}

\begin{keyconceptbox}[title=Convolution in Control Systems]
The output $y(t)$ of a linear time-invariant (LTI) system with impulse response $h(t)$ to input $u(t)$ is:
\begin{equation}
y(t) = (h \conv u)(t) = \int_0^t h(\tau) u(t - \tau) \, \dd\tau
\end{equation}
\end{keyconceptbox}

% ============================================================================
